\section{Polynomial curves, tangents, and a quantitative implicit theorem}
\label{sec:tangency}

The source article's informal ``just touches'' description does not determine
a tangent. For example, the tangent $y=0$ to $y=x^3$ at the origin crosses
the graph. We use a finite quadratic remainder as the definition of first-order
contact. This supplies an explicit \ipra\ interpretation of tangent lines and
planes, and also proves existence of the relevant local curve or surface when
a partial derivative is separated from zero.

\subsection{Finite derivatives and remainder estimates}

Let $F(X)=\sum_\alpha c_\alpha X^\alpha$ be a rational polynomial in a
fixed number $d$ of variables. The exponent vector is $\alpha$, its total
degree is $|\alpha|$, and coefficients and exponents are finite coded lists.
Define formal partial derivatives term by term, and write $\nabla F$ for
their vector. For $M\ge1$ define the rational bounds
\begin{align*}
 H_F(M)&=\sum_\alpha |c_\alpha|(M+1)^{|\alpha|},\\
 B_F(M)&=1+\sum_\alpha |c_\alpha|\,|\alpha|(M+1)^{|\alpha|},\\
 L_F(M)&=1+\sum_\alpha |c_\alpha|\,|\alpha|^2(M+1)^{|\alpha|}.
\end{align*}
All are uniformly primitive recursive in the polynomial code and $M$.

\begin{theorem}[A polynomial Taylor estimate without a limit theorem]
\label{tan:taylor}
If $\|p\|_\infty\le M$ and $\|h\|_\infty\le1$, then
\[
 F(p+h)=F(p)+\nabla F(p)\cdot h+R(p,h),\qquad
 |R(p,h)|\le H_F(M)\|h\|_\infty^2.
\]
The equality is a finite polynomial identity and the bound is a rational
inequality. Both are provable in \ipra, uniformly for coded polynomials.
\end{theorem}
\begin{proof}
For one monomial expand each factor by the binomial formula:
\[
 (p+h)^\alpha=\sum_{\beta\le\alpha}
 {\alpha\choose\beta}p^{\alpha-\beta}h^\beta.
\]
The term $|\beta|=0$ is $p^\alpha$ and the terms $|\beta|=1$ are
the formal derivatives times $h_i$. In every remaining term,
$|h^\beta|\le\|h\|_\infty^2$, since $\|h\|_\infty\le1$.
The sum of the remaining coefficient bounds is at most the sum of all
of them,
\[
 \sum_{\beta\le\alpha}{\alpha\choose\beta}M^{|\alpha|-|\beta|}
       =(M+1)^{|\alpha|}.
\]
Multiply by $|c_\alpha|$ and sum over the finite coefficient list.
The binomial identity and list sums are proved by the finite induction
rule of Lemma~\ref{found:loops}. Transfer to fixed PR names uses
Proposition~\ref{found:transfer} and the same explicit bounds.
\end{proof}

\begin{corollary}[Tangent lines and planes to polynomial graphs]
\label{tan:graphs}
For a polynomial $f:\Q^m\to\Q$, $m=1$ or $2$, the graph at
$(p,f(p))$ has the affine first-order approximation
\[
 z=f(p)+\nabla f(p)\cdot(x-p).
\]
Its vertical error on $\|x-p\|_\infty\le1$ is at most
$H_f(M)\|x-p\|_\infty^2$. A normal vector is
$(-\nabla f(p),1)$. This affine approximation is unique among affine
functions whose residual is bounded by a constant times
$\|x-p\|_\infty^2$ near $p$.
\end{corollary}
\begin{proof}
The remainder estimate is Theorem~\ref{tan:taylor}. A displacement
$(h,\nabla f(p)\cdot h)$ in the affine graph has zero dot product
with the stated normal. Putting $h=0$ forces both affine approximations to agree at $p$.
Subtract them and substitute $h=te_i$. Their coefficient difference
$a_i$ satisfies $|a_i|\,|t|\le C t^2$ for all sufficiently small
nonzero rational $t$. If a fixed rational observation of $a_i$ had a
positive gap from zero, a dyadic $t$ smaller than that gap divided by
$2(C+1)$ would contradict the estimate. Therefore the coefficient
names are equal at every precision. This proof uses an explicit small
rational $t$, not an assumption about infinitesimals.
\end{proof}

\begin{theorem}[Polynomial parametrized curves and tangent transport]
\label{tan:parametric}
Let $\gamma(t)\in\Q[t]^d$, $d=2$ or $3$. At a rational parameter $t_0$ the vector
$v=\gamma'(t_0)$ satisfies
\[
 \|\gamma(t_0+h)-\gamma(t_0)-hv\|_\infty\le Hh^2
 \quad(|h|\le1),
\]
where $H$ is the maximum of the component Taylor bounds. If $v\ne0$,
the tangent is $\gamma(t_0)+\R v$ in the represented-family reading.
Under a polynomial reparametrization $t=\sigma(u)$ with
$\sigma'(u_0)\ne0$ and $\sigma(u_0)=t_0$, the tangent line is
unchanged. If $F(\gamma(t))$ is the zero polynomial, then
$\nabla F(\gamma(t_0))\cdot v=0$.
\end{theorem}
\begin{proof}
Apply Theorem~\ref{tan:taylor} to each component. The finite product
rule proves the formal chain rule, first on a monomial by distributing
its one-factor derivative, then on a finite sum. It gives
$(\gamma\circ\sigma)'(u_0)=v\sigma'(u_0)$, a nonzero scalar
multiple of $v$, and gives the last dot product by differentiating
the zero polynomial. A coefficientwise zero polynomial has all its
formal derivatives zero, so this last argument is exact. The regularity
conditions are decidable for rational input; for named input they
require the stated nonzero gaps before normalization or inversion.
No claim of a unique tangent is inferred from a vanishing velocity.
\end{proof}

\begin{lemma}[Explicit divided-difference bounds]
\label{tan:divided}
Suppose $\|p\|_\infty\le M$ and all points under consideration are
within distance $r\le1$ of $p$ in the maximum norm. If only coordinate
$i$ changes from $u$ to $v$, write the exact identity
\[
 F(\ldots,v,\ldots)-F(\ldots,u,\ldots)
          =(v-u)D_iF(\ldots;u,v).
\]
Here $D_iF$ is a polynomial, including at $u=v$. It obeys
\[
 |D_iF-\partial_iF(p)|\le L_F(M)r.
\]
More generally $|F(q)-F(q')|\le B_F(M)\|q-q'\|_\infty$
on that box. These estimates require neither integration nor a
mean-value theorem.
\end{lemma}
\begin{proof}
For a power use
\[
 \frac{v^k-u^k}{v-u}=\sum_{j=0}^{k-1}v^{k-1-j}u^j,
\]
interpreted as a polynomial identity even when $u=v$. For a monomial
of total degree $a$, its divided difference has at most $a$ products,
each with $a-1$ factors. Compare each factor to the corresponding
coordinate of $p$ and telescope. Each factor difference is at most
$r$, every remaining factor has magnitude at most $M+1$, so the
total error is bounded by
$a(a-1)(M+1)^{\max(a-2,0)}r$. The definition of $L_F$ is a larger
bound after multiplying by coefficients and summing. For the second
inequality telescope directly in all factors of each monomial:
its error is at most $a(M+1)^{\max(a-1,0)}\|q-q'\|_\infty$.
Summing is bounded by $B_F(M)$. This last estimate bounds the whole
coordinate change, so no missing factor of the dimension is needed.
\end{proof}

\subsection{A local implicit curve or surface with a PR constructor}

\begin{theorem}[Quantitative polynomial implicit function]
\label{tan:implicit}
Let $F(x,y)$ be a rational polynomial, where $x\in\Q^m$,
$m=1$ or $2$, and let $p=(x_0,y_0)$ be a rational point with
\[
 F(p)=0,\qquad g=\partial_yF(p)>0.
\]
Choose an integer $M\ge1$ with $\|p\|_\infty\le M$ and put
\[
 B=B_F(M),\quad L=L_F(M),\quad
 r=\min\left(1,\frac{g}{4(L+1)}\right),\quad
 a=\frac{gr}{4(B+g+1)}.
\]
For every rational $x$ with $\|x-x_0\|_\infty\le a$, there is
a PR name $Y(x,n)$ for a unique $y(x)\in[y_0-r,y_0+r]$ satisfying
$F(x,y(x))=0$. It is obtained by exactly $n+2$ rational bisections
of $[y_0-r,y_0+r]$ and taking the midpoint. Its interval width is
at most $2^{-n-1}$. In addition,
\[
 |y(x)-y(x')|\le\frac{2B}{g}\|x-x'\|_\infty
\]
for rational $x,x'$ in this box. Every displayed conclusion is a
fixed-family first-order arithmetic statement. If the nonzero partial
derivative is negative, replace $F$ by $-F$; if another partial is
nonzero, permute coordinates first.
\end{theorem}
\begin{proof}
Since $a<r$, all points used are in the radius-$r$ box. By
Lemma~\ref{tan:divided}, every divided difference in the $y$ direction
differs from $g$ by at most $Lr\le g/4$, so is at least $g/2$.
Thus for rational $u<v$ in the vertical interval,
\[
 F(x,v)-F(x,u)\ge(g/2)(v-u).
\]
This is a strict monotonicity inequality with an explicit positive
coefficient. Also
\[
 |F(x,y_0)|\le B\|x-x_0\|_\infty\le Ba<gr/4.
\]
Moving vertically by $r$ changes $F$ by at least $gr/2$.
Consequently $F(x,y_0-r)<0<F(x,y_0+r)$, with rational sign margins.

Start with those two endpoints. At a rational midpoint test the exact
sign of $F$: if it is nonpositive, retain the upper half, and otherwise
retain the lower half. The invariant is
\[
 F(x,l_j)\le0\le F(x,u_j),\qquad
 [l_{j+1},u_{j+1}]\subseteq[l_j,u_j],\qquad
 u_j-l_j=2r\,2^{-j}.
\]
It consists of rational comparisons and is proved by quantifier-free
induction. If a midpoint is an exact root, the chosen interval still
has that root as an endpoint and preserves the invariant. At $j=n+2$
the width is at most $2^{-n-1}$. Nested intervals give the name condition
and a named point in the original vertical interval.

For the midpoint $q$, the endpoint signs and the Lipschitz estimate give
$|F(x,q)|\le B(u_j-l_j)$. To obtain residual at most $2^{-k}$ take
$n=k+\ell(\lceil B\rceil+1)+2$. Thus the name satisfies the polynomial
equation at all precisions. This establishes existence directly by a
bounded procedure, not by an appeal to an arbitrary intermediate-value
principle.

If two named roots $y,y'$ are in the interval, the divided-difference
identity yields $|F(x,y)-F(x,y')|\ge(g/2)|y-y'|$ in the finite
observation interpretation, and hence $y=y'$. This inequality transfers
because its divided-difference coefficient has the uniform positive
bound already proved; no sign test on $y-y'$ is performed. Finally,
\[
 (g/2)|y(x)-y(x')|
 \le |F(x,y(x))-F(x,y(x'))|
 =|F(x',y(x'))-F(x,y(x'))|
 \le B\|x-x'\|_\infty.
\]
It proves the Lipschitz statement and completes the primitive-recursive
construction. All parameters in the number of bisections and in the
precision request are explicit rational or natural data.
\end{proof}

\begin{corollary}[The tangent of the implicit branch, with its error]
\label{tan:implicit-tangent}
Under Theorem~\ref{tan:implicit}, set $K=\max(1,2B/g)$ and
$f=\nabla_xF(p)$. Then
\[
 \left|y(x)-y_0+\frac{f\cdot(x-x_0)}{g}\right|
 \le \frac{H_F(M)K^2}{g}\|x-x_0\|_\infty^2.
\]
In particular the tangent line or plane is exactly
$\nabla F(p)\cdot(X-p)=0$, and the gradient is a normal vector.
\end{corollary}
\begin{proof}
The Lipschitz estimate with $x'=x_0$ bounds the vertical displacement by
$(2B/g)\|x-x_0\|_\infty$. Insert
$h=(x-x_0,y(x)-y_0)$ into Theorem~\ref{tan:taylor}. Both polynomial
values vanish, so
$0=f\cdot(x-x_0)+g(y(x)-y_0)+R$, with
$|R|\le H_F(M)K^2\|x-x_0\|_\infty^2$. Divide by the positive
rational $g$. The linear equation is the kernel of the nonzero
gradient covector, hence a line when $m=1$ and a plane when $m=2$.
Uniqueness follows by the coefficient argument of
Corollary~\ref{tan:graphs}.
\end{proof}

\begin{proposition}[Regularity is an essential hypothesis]
\label{tan:singular}
The equations $y^2-x^2=0$ and $y^2-x^3=0$ both have vanishing gradient
at the origin, but the first has two distinct rational line branches,
while the parametrized cusp $(t^2,t^3)$ of the second has tangent
direction the horizontal axis in the limiting-secant sense. Thus a
zero gradient does not by itself select a tangent, or certify its
nonexistence.
\end{proposition}
\begin{proof}
The first polynomial is $(y-x)(y+x)$; the two displayed line
parametrizations have respective directions $(1,1)$ and $(1,-1)$.
They are distinct rational directions. For the second, direct
substitution gives $(t^3)^2=(t^2)^3$. For $t\ne0$, its secant vector
from the origin is $t^2(1,t)$, and the slope is $t$. This tends to zero
with the explicit bound $|t|$. At $t=0$ the parametrization's velocity
vanishes, so Theorem~\ref{tan:parametric}'s regular-velocity conclusion
was not applicable. These computations establish both examples without
making any general decision claim about singularities.
\end{proof}

\subsection{Arc length as an explicitly convergent geometric observable}

\begin{theorem}[Constructive arc length of a polynomial curve]
\label{tan:arclength}
Let $\gamma(t)=\sum_j c_jt^j\in\Q[t]^d$, $d=2$ or $3$, on a
rational interval $[a,b]$, $a<b$. Put $D=b-a$ and choose
$M\ge\max(1,|a|,|b|)$. With component coefficients $c_{ij}$ set
\[
 V=\sum_{i=1}^{d}\sum_{j\ge1}|c_{ij}|jM^{j-1},\qquad
 K=\sum_{i=1}^{d}\sum_{j\ge2}|c_{ij}|j(j-1)M^{j-2}.
\]
Let $L_N$ be the sum of the Euclidean chord lengths over the uniform
$N$-interval partition. The dyadic sequence $L_{2^k}$ is nondecreasing,
bounded by $VD$, and has the explicit Cauchy bound
\[
 0\le L_{2^{k+s}}-L_{2^k}\le 2KD^2\,2^{-k}.
\]
It therefore defines a PR name for arc length. This name is at least
the chord sum of every finite rational partition and is approximated by
those sums as their mesh tends to zero with the displayed modulus.
\end{theorem}
\begin{proof}
For any subinterval $[u,v]\subseteq[a,b]$ the polynomial divided
difference proves
\[
 \left\|\frac{\gamma(v)-\gamma(u)}{v-u}\right\|\le V.
\]
The same factorwise comparison to $\gamma'(t_0)$, when $u,v,t_0$
are in an interval of width $h$, gives
\[
 \|\gamma(v)-\gamma(u)-(v-u)\gamma'(t_0)\|
                  \le K h(v-u).
\]
The Euclidean norm is bounded by the sum of coordinate absolute values,
so the displayed $V,K$ suffice. Both estimates follow by finite
telescoping, exactly as in Lemma~\ref{tan:divided}.

Fix a coarse interval of width $h$ and choose its left endpoint as
$t_0$. Each refinement chord differs in length from
$(v-u)\|\gamma'(t_0)\|$ by at most $Kh(v-u)$, using the metric
inequality $|\|A\|-\|B\||\le\|A-B\|$. Summing the refined
chords gives error at most $Kh^2$ from $h\|\gamma'(t_0)\|$.
The coarse chord has the same bound. Thus refinement increases the
chord sum by at most $2Kh^2$; increase is nonnegative by the triangle
inequality. Over $N$ equal coarse intervals this is $2KD^2/N$,
independently of how finely they are subdivided. The bound $VD$ follows
by summing the first inequality.

To make an actual rational output at precision $n$, let
$C=\lceil2KD^2\rceil+1$ and choose
$k=n+\ell(C)+3$, $N=2^k$. Evaluate the rational coordinates at
the partition points exactly. Each chord has a nonnegative rational
squared length. Approximate its square root at precision
$m=n+k+3$ using Theorem~\ref{found:sqrt}. The sum of the $N$ errors
is at most $2^{-n-3}$; the refinement tail is at most $2^{-n-3}$.
Thus the rational output has a uniform PR Cauchy modulus with slack.

For a nonuniform partition with mesh $h$, the same argument gives a
refinement error at most $2K\sum h_i^2\le2KD h$. Take its common
refinement with a fine dyadic partition. This shows its chord sum is
bounded above by the limiting dyadic name, and that its deficit is
at most $2KD h$. Finite partition merging is bounded list arithmetic.
This characterizes the usual polygonal arc length for this fixed PR
family using only natural-number codes of finite partitions, with no
unrestricted real supremum principle inside the theory.
\end{proof}

The theorem includes curves with zero velocity and self-intersections;
it measures parametrized length, with multiplicity when the curve retraces
itself. It is not a theorem about arc length of every continuous curve.
The finite observation modulus, and the distinction between a curve's
parametrization and its image, are part of the proved statement.
